% A.104 reserved by the root.  Do not add an active input before audit.
\begin{statement}{A.104}{Proved}
{The four simultaneous projection classes and their detection at three}
Fix $\lambda\in\Q\setminus\{0,1\}$ and its point
section $s_\lambda$.  Put
\[
 \begin{split}
  C_H&=\ker(\Pi\longrightarrow H),\qquad
  E_J=E_U/(C_H\cap\Gamma_2\Pi),\\
  T&=\Pi^{\rm ab}=\widehat\Z(1)^2,\qquad
  M=H^{\rm ab},\qquad K=\ker(T\longrightarrow M).
 \end{split}
\tag{A.104.1}
\]
For $k=\Q$ or $\Q_3$, let $\mathcal F_k(\lambda)$
be the set of geometric conjugacy classes of continuous
sections of $(E_J)_k\to G_k$ whose projections to
$(E_H)_k$ and $(E_1)_k$ have the same geometric
conjugacy classes as the respective projections of
$s_\lambda$.  Here $E_1=E_U/\Gamma_2\Pi$, and the
local extensions are the pullbacks along the chosen
map $G_{\Q_3}\to G_\Q$.

Both sets have exactly four elements.  Relative to
the point splitting they have the exact description
\[
 \mathcal F_k(\lambda)
      \simeq M^{G_k}
      =\{(4\epsilon_0,4\epsilon_1,0):
                       \epsilon_0,\epsilon_1\in\{0,1\}\},
\tag{A.104.2}
\]
using A.56's coordinates
$M=(\Z/8(1))^2\oplus\Z/3(1)$.
Restriction from $\Q$ to $\Q_3$ is a bijection
between these four-element sets, identified with
the identity on (A.104.2).
Only the zero element has the local class of an
ordinary point or cusp at three.

For $E_\lambda$ good at three, the three nonzero
classes already have no ordinary or cuspidal source
in the explicit
geometric quotient
\[
       J_{32,16}=
       (H_2\bmod32)
          \mathbin{\times}_{(\Z/8(1))^2}
                            (\Z/16(1))^2,
                 \qquad |J_{32,16}|=8192.
\tag{A.104.3}
\]
The arithmetic quotient uses the action induced
by $s_\lambda$; the displayed fibre product and its
map from $\Pi$ are the ones induced by A.56.
This includes every rational Legendre parameter
with $j=1728$.
If $E_\lambda$ is potentially multiplicative at
three, put $e=-v_3(j(E_\lambda))/2$ and
$N=5+v_2(e)$.  Replacing $H_2\bmod32$ in
(A.104.3) by $H_2\bmod2^N$ gives an explicit
quotient of order $2^{3N-2}$ that distinguishes
the four local classes from one another.

Thus the two projection classes by themselves
admit three additional quotient classes, with
exact relative-sign parameters (A.104.2), but
none of the additional classes matches locally
at three.  The statement classifies this specified
simultaneous projection fibre, not the whole
fibre over the $H$ projection alone.
No local section theorem is used.
\end{statement}
\begin{proof}
\emph{The actual fibre-product extension and its cocycles.}
A.56 identifies the geometric group and the full
arithmetic extension as
\[
          J=H\mathbin{\times}_M T,
                     \qquad E_J=J\rtimes G_\Q,
\tag{A.104.4}
\]
where the splitting is $s_\lambda$.
The kernel of $J\to H$ is the group $K$ in
(A.104.1), embedded as $(1,k)$.
It is central in $J$, since the second factor
$T$ is abelian, and its $G_\Q$-action is the
restriction of the cyclotomic action on $T$.
All of these maps are continuous homomorphisms
of the displayed profinite extensions.

First fix the $H$ projection as a literal
section, rather than just a section class.
A section with this projection has the unique
form $z\cdot s_\lambda$, where
$z:G_k\to K$ is a continuous one-cocycle.
Conjugation by $K$ changes it by a $K$-coboundary.
Its abelian projection has the point's class
exactly when the image of $[z]$ in
$H^1(G_k,T)$ is zero.
There is an exact injection
\[
 \begin{split}
 \delta_k:M^{G_k}&\longrightarrow H^1(G_k,K),\\
 v&\longmapsto[\,g\longmapsto g t_v-t_v\,],
                  \qquad \eta(t_v)=v,
 \end{split}
\tag{A.104.5}
\]
whose image is precisely this kernel;
$\eta:T\to M$ is the surjection of A.56.
Here $T^{G_k}=0$ for both fields: an invariant
in each $\Z_\ell(1)$ factor would give a
compatible system of $\ell$-power roots of
unity in $k$, whereas the roots of unity in
either field form a finite group.

For completeness, exactness in (A.104.5) is
an assertion about continuous cocycles, with
no passage through uncomputed derived limits.
The finite group $M$ has chosen lifts in $T$,
so each displayed cocycle is continuous and
lies in $K$ when $v$ is invariant.
Changing $t_v$ by an element of $K$ changes
it by a $K$-coboundary.
If it were a $K$-coboundary, subtracting that
element would leave an element of $T^{G_k}$,
so $v=0$.  Conversely, a $K$-cocycle that
is a $T$-coboundary has exactly this form;
its chosen $t$ has invariant image in $M$.
This proves both the injection and its image.

We must also account for forgetting the literal
$H$ comparison.  Every geometric $H$-conjugator
lifts through $J\to H$, by the surjectivity
proved in A.56, so every class in
$\mathcal F_k(\lambda)$ has a representative
of the preceding form.
Two such representatives can be conjugate
only by an element $(h,t)\in J$ for which
$h\in H^{G_k}$ centralizes the point section.
Under the fibre-product relation one has
$\eta(t)=h^{\rm ab}$, and the actual
conjugation formula replaces $z(g)$ by
\[
                        z(g)+t-g t.
\tag{A.104.6}
\]
Conjugation acts trivially on the central
group $K$ itself.
Consequently forgetting the literal comparison
quotients the parameters in (A.104.5) exactly
by the image of
\[
                   H^{G_k}\longrightarrow M^{G_k},
                                  \qquad h\longmapsto h^{\rm ab}.
\tag{A.104.7}
\]
We now prove that this image is zero for both
fields, including every special parameter.

\emph{The two-adic centralizer at the fixed physical prime three.}
The actual arithmetic representation map
$E_H\to\operatorname{GL}_2(\widehat\Z)\times G_\Q$
is injective: an element in its kernel is
geometric, and the geometric representation
has faithful image $H$ by definition.
Thus $H^{G_{\Q_3}}$ is exactly the geometric
centralizer of the point's local Tate system.

At three a rational Legendre curve is either
good or potentially multiplicative.
If $\lambda$ and $1-\lambda$ are both
three-adic units, the integral Legendre
model has unit discriminant.  Otherwise
the formula
\[
             j(E_\lambda)=
              \frac{256(\lambda^2-\lambda+1)^3}
                   {\lambda^2(1-\lambda)^2}
\tag{A.104.8}
\]
has negative three-adic valuation, giving
the second alternative.

In the potentially multiplicative case,
pass to a finite extension on which the
curve is a split Tate curve.
Its rational two-adic cyclotomic line is
unique, and its extension class has
nonzero ordinary Tate valuation.
The centralizer computation (A.78.7)
then makes every commuting rational matrix
scalar: equality of the two diagonal
entries follows from the nonzero extension
class, and its upper entry vanishes since
the cyclotomic character is nontrivial.
The original local centralizer is a
subgroup of this one.
An integral scalar in $H_2$ has square
one by its determinant; of $I,-I$,
only $I$ has diagonal entries one modulo
four.  Hence its image in $H_2^{\rm ab}$
is zero.

In the good case, both units must reduce
to two modulo three, so the reduction of
$E_\lambda$ is literally
$Y^2=X(X-1)(X-2)$ over $\F_3$.
It has four rational points.
Arithmetic Frobenius $F_3$ on $T_2$ therefore
has trace zero and determinant three.
The exact mod-four formula of A.16, with
$\chi_4(\operatorname{Frob}_3)=3$ and
both unit Kummer characters nontrivial,
gives
\[
           F_3\equiv\begin{pmatrix}3&2\\2&1\end{pmatrix}
                 \pmod4,
       \qquad B=(F_3-I)/2\in M_2(\Z_2),
       \qquad B^2+B+I=0.
\tag{A.104.9}
\]
In particular $\operatorname{tr}B=-1$,
$\det B=1$, and the characteristic
polynomial of $B\bmod2$ is the irreducible
polynomial $X^2+X+1$.
For any primitive vector $v$, the pair
$v,Bv$ is a $\Z_2$-basis: its reduction
is independent because $B\bmod2$ has no
invariant line.  A commuting matrix is
therefore determined by its value on $v$
and is exactly $aI+bB$, with $a,b\in\Z_2$.
This proves the integral centralizer
description, rather than only its
rational counterpart.

Suppose $h=aI+bB\in H_2$.
Both off-diagonal entries of $B$ are odd,
and its two diagonal entries have distinct
parities by (A.104.9).
The off-diagonal and diagonal conditions
defining $H_2$ force
\[
        a=1+4u,\qquad b=4v\quad(u,v\in\Z_2),
              \qquad
        1=\det(h)=a^2-ab+b^2.
\tag{A.104.10}
\]
Reducing the last equality modulo sixteen
gives $v\equiv2u\pmod4$.
Consequently
\[
 \begin{split}
       h&\equiv I+4uF_3\equiv c^u\pmod{16},\\
       c&=\begin{pmatrix}13&8\\8&5\end{pmatrix}
                    =[U(2),L(-2)].
 \end{split}
\tag{A.104.11}
\]
The equality $c^u\equiv I+4uF_3$ follows
by reducing $u$ modulo four and observing
that the square of $4F_3$ is zero modulo
sixteen.
A.56 proves the full inverse-limit
identity
\[
       [H_2,H_2]
            =\{h\in H_2:h\bmod16\in\langle c\rangle\};
\tag{A.104.12}
\]
in particular it includes the entire
level-sixteen kernel, not merely its
first finite approximation.
Equations (A.104.11)--(A.104.12) show that
the full two-adic centralizer has zero
image in $H_2^{\rm ab}$.

For both $k=\Q$ and $\Q_3$ the invariant
group of $M$ is
\[
       M^{G_k}=\mu_8(k)^2\oplus\mu_3(k)
                         =\{\pm1\}^2\oplus\{1\}.
\tag{A.104.13}
\]
At three, a fourth root of unity would
make $-1$ a square modulo three, and a
primitive cube root would make $-3$
a square in $\Q_3$; neither is possible.
These observations give exactly the
roots used in (A.104.13).
Thus there is no three-primary invariant
to receive a centralizer image, and
the two-primary image has just been
proved zero.
This proves (A.104.7) is zero over
$\Q_3$.  A global centralizer restricts
to the local one, so its image is zero
over $\Q$ as well.
No global endomorphism theorem is needed
for this implication.

\emph{The four explicit representatives and localization.}
Choose once a compatible Tate basis for
the rank-one cyclotomic module.
For $\epsilon=(\epsilon_0,\epsilon_1)\in\{0,1\}^2$ put
\[
 \begin{split}
       v_\epsilon&=(4\epsilon_0,4\epsilon_1,0)\in M,\\
       t_\epsilon&=(12\epsilon_0,12\epsilon_1)\in T,\\
       z_\epsilon(g)&=(\chi(g)-1)t_\epsilon\in K.
 \end{split}
\tag{A.104.14}
\]
The equality $\eta(t_\epsilon)=v_\epsilon$
uses both mod-eight coordinates and the
zero sum modulo three.
Each displayed map is a continuous
cocycle, so pairing its abelian section
with the literal $H$ point section
defines a section of $E_J$.
Its abelian projection is the point
section conjugated by $-t_\epsilon$,
as the semidirect-product formula
$(-t,1)(0,g)(t,1)=(g t-t,g)$ shows.
The element $(1,-t_\epsilon)$ need not
belong to $J$; its failure to belong
is precisely the nonzero invariant
$v_\epsilon$.
The injection (A.104.5) and the zero
centralizer image prove that the four
resulting geometric section classes
are distinct and exhaustive.

Restriction sends (A.104.14) to the
same formula over $\Q_3$.
Equation (A.104.13) identifies its
invariant group with the same four
elements, so restriction is a bijection.
For a local ordinary point $x$ whose
abelian projection equals that of
$\lambda$, injectivity of the full
local Kummer map gives $x=\lambda$.
For the cusp packets, the regular
coordinates $1-t,t,t/(t-1)$ have
Kummer class one at cusps $0,1,\infty$,
respectively.  The corresponding point
values $1-\lambda,\lambda,-\lambda/(1-\lambda)$
are never one for $\lambda\ne0,1$.
Local Kummer injectivity therefore
excludes every cusp.
Only the zero class of
$\mathcal F_{\Q_3}(\lambda)$ has an
ordinary or cuspidal source, proving
the asserted local obstruction for
all three other global classes.

\emph{An explicit finite detector in the good-at-three case.}
Project (A.104.4) to its two-primary
simultaneous factor, and then reduce
$H_2$ modulo $32$ and $T_2$ modulo $16$.
A.56's geometric surjectivity and
abelianization formula give the exact
quotient (A.104.3), with its point-induced
arithmetic action.  The first factor
has $2048$ elements, and the kernel
of $(\Z/16)^2\to(\Z/8)^2$ has four,
proving its order.

Let $h\in H_2\bmod32$ commute with
the local Frobenius modulo $32$.
Since $F_3=I+2B$, its reduction modulo
$16$ commutes with $B\bmod16$.
The cyclic-basis argument following
(A.104.9) works over $\Z/16$ as well.
It again gives $h=aI+bB\pmod{16}$;
the $H_2$ conditions and determinant
modulo sixteen give exactly
(A.104.10)--(A.104.11) at that level.
Thus $h$ has zero image in
$(\Z/8)^2$.

A conjugator in (A.104.3) that leaves
the $H$ point projection fixed must
therefore have its $(\Z/16)^2$
coordinate divisible by eight.
Its abelian coboundary at arithmetic
Frobenius three is
$(3-1)t=0\pmod{16}$.
But the representative (A.104.14)
has value
\[
                z_\epsilon(\operatorname{Frob}_3)
                  =(8\epsilon_0,8\epsilon_1)\pmod{16}.
\tag{A.104.15}
\]
For nonzero $\epsilon$ this cannot be
removed by such a conjugator.
The finite quotient thus already
distinguishes all four local classes.

It also excludes every other ordinary point
and every cusp as a possible source of a
nonzero class in this finite quotient.
First take an ordinary point $x\in\Q_3$
whose Legendre curve has good reduction.
Then $x,1-x$ are units, so $x\equiv\lambda\equiv2$
modulo three.  The actual residue-disc map is
\[
 \operatorname{Spec}\Z_3[[z]][1/3]\longrightarrow U_{\Q_3},
                           \qquad t=\lambda+z.
\]
Both $t$ and $1-t$ are units in its source
ring.  Evaluation at $z=0$ and at
$z=x-\lambda$ is defined by convergent
integral power series and gives the two
indicated points.
The geometric pro-two group of this
unpunctured scheme is trivial by the
normalization proof of A.49, now with
only the vertical boundary divisor.
More explicitly, a finite geometric
two-group cover descends over finite
constants; its finite normalization
over the excellent ring $\mathcal O_F[[z]]$
can ramify only at the vertical divisor,
where it is tame.  A constant extension
adjoining a suitable root of a uniformizer
kills that ramification.  Purity makes the
normalization etale, and an unramified
constant extension splits it by the
henselian local equivalence.
Over algebraic constants it is therefore
trivial.  The corresponding arithmetic
extension has zero geometric kernel,
so the two point sections have the same
class in every geometric pro-two quotient,
in particular (A.104.3).

For an ordinary point with potentially
multiplicative reduction, or for any
of the three cusp packets, the actual
coefficient-two Tate comparison of A.63
has a quadratic scalar character $\xi$
and diagonal characters $\xi\chi_2,\xi$.
If $\xi$ is ramified at three, an inertia
element with $\xi=-1$ has trace $-2$
modulo $32$, whereas the good point's
inertia has trace $2$.
If $\xi$ is unramified, Frobenius has
trace $4$ or $-4$ modulo $32$, whereas
the good point has trace zero.
These contradictions exclude all such
sources already in the $H_2\bmod32$
projection.  Every ordinary Legendre
point at three is in one of these two
reduction cases, by (A.104.8).
Together with the residue-disc comparison,
this proves the stated finite point-or-cusp
exclusion.

Finally the exact polynomial identity
\[
 \begin{split}
  &256(\lambda^2-\lambda+1)^3
       -1728\lambda^2(1-\lambda)^2\\
       &\hspace{2em}=64(\lambda+1)^2(\lambda-2)^2(2\lambda-1)^2
 \end{split}
\tag{A.104.16}
\]
shows that the rational parameters
with $j=1728$ are $-1,2,1/2$.
All reduce to two modulo three and
have good reduction there, so the
same finite detector treats them.
This closes the special case without
an assumption about a local CM
centralizer.

\emph{The explicit finite bound in the potentially multiplicative case.}
Suppose finally that $E_\lambda$ is
potentially multiplicative at three.
Choose the actual small cusp coordinate
$w=\lambda$, $1-\lambda$, or $1/\lambda$,
according to the contact label.
Its valuation is the positive integer
$e=-v_3(j)/2$.
The inverse-nome expression of A.63 gives
\[
         q_3=\frac{w^2}{256}F(w)
                 =3^{2e}u_0,
                       \qquad u_0\in1+3\Z_3.
\tag{A.104.17}
\]
The unit assertion follows because the
square unit $w^2/3^{2e}$ and $256$ both
reduce to one modulo three, and
$F(w)\equiv1\pmod3$.
Every principal unit has unique compatible
two-power roots in $1+3\Z_3$, by Hensel's
lemma.  Choosing these actual roots makes
the Kummer cocycle of $q_3$ exactly
$2e$ times that of $3$.

In the tame extension
$\Q_3(\mu_{2^\infty},3^{1/2^\infty})$
choose $\tau$ fixing the roots of unity
and multiplying each $2^n$th root of
$3$ by the corresponding root of unity.
Choose Frobenius $\varphi$ fixing these
roots of $3$ and raising roots of unity
to their third powers.
These compatible automorphisms exist
at each finite level: the cyclotomic
extension is unramified, the Eisenstein
root extension is totally ramified,
and their compositum has the indicated
root maps.  They extend to elements
of $G_{\Q_3}$ by restriction surjectivity.
After cancelling the quadratic scalar
Tate twist, the actual integral matrices
in Tate coordinates are
\[
          \rho(\tau)=\begin{pmatrix}1&2e\\0&1\end{pmatrix},
          \qquad
          \rho(\varphi)=\begin{pmatrix}3&0\\0&1\end{pmatrix}.
\tag{A.104.18}
\]
Scalar quadratic twists do not affect
their centralizer equations.

Let $N=5+v_2(e)$.
A matrix commuting with the second
matrix modulo $2^N$ has off-diagonal
entries divisible by $2^{N-1}$.
Commutation with the first makes its
diagonal difference divisible by
$2^{N-1-v_2(e)}=16$.
It is therefore scalar modulo sixteen.
Scalarity is preserved on returning
through the actual integral change
from Tate coordinates to the original
Legendre basis.
For an element of $H_2$, its scalar
residue is one modulo four and its
square is one modulo sixteen.
The only possibilities are $I$ and
$9I=c^2$, both in the derived image
of A.56.  Thus the centralizer modulo
$2^N$ again has zero image in
$(\Z/8)^2$.
The same Frobenius calculation
(A.104.15) excludes each nonzero
relative-sign parameter in
\[
       J_{2^N,16}
        =(H_2\bmod2^N)
              \mathbin{\times}_{(\Z/8(1))^2}
                               (\Z/16(1))^2,
          \qquad |J_{2^N,16}|=2^{3N-2}.
\tag{A.104.19}
\]
Indeed $|H_2\bmod2^N|=2^{3N-4}$
and the other projection has kernel
of order four.
This proves the asserted explicit
finite distinction in the remaining
reduction case.
\end{proof}

The certificate
\path{code/legendre_simultaneous_fibre_certificate.py},
with its retained output, verifies the finite
quotient, all $128$ possible Frobenius matrices
modulo $32$ satisfying (A.104.9), and all $4096$
commuting pairs with $H_2\bmod32$.
Their abelian images are all zero.
It also checks the four invariant and cocycle
formulas and the displayed $j$ factorization.
The full continuous-cohomology classification
and the potentially multiplicative case are
proved above.
